\chapter{Literature review}

\section{Introduction}
The resilience of transport networks has become a critical focus in infrastructure planning, particularly in regions prone to disruptions such as natural disasters, extreme weather events, and structural failures. A resilient transport network ensures continued connectivity, minimizes delays, and provides alternative routes when disruptions occur. Resilience in transport networks refers to the system's ability to absorb, adapt to, and recover from disruptions while maintaining essential functions and services \cite{holling1973}. It encompasses network robustness, redundancy, and adaptive capacity, ensuring that transport infrastructure can continue functioning under various stress conditions \cite{taylor2017}.

Closely related to resilience is the concept of redundancy, which refers to the presence of alternative routes or backup infrastructure that can sustain functionality when primary links fail. A network with high redundancy can reallocate traffic through alternative paths, reducing the impact of disruptions and preventing system-wide failures \cite{zhang2018}. Redundancy is particularly important in interregional transport networks, where disruptions to major highways, bridges, or rail lines can have widespread economic and social consequences.

This chapter reviews key studies and methodologies used to evaluate transport network resilience, with a particular focus on graph theory indicators and GIS-based spatial analysis. Graph theory provides a quantitative framework for analyzing network connectivity, redundancy, and structural robustness, while Geographic Information Systems (GIS) facilitate spatial modeling and visualization of disruptions. The combination of these methodologies allows researchers to identify critical links and nodes, assess network redundancy, and propose infrastructure improvements aimed at enhancing resilience.

The literature review is structured into several sections. The first section discusses the role of graph theory indicators in assessing transport network performance, including key metrics such as beta index, gamma index, clustering coefficient, and betweenness centrality. These indicators provide insights into network complexity, connectivity, and the impact of disruptions on overall system functionality. The next section focuses on seismic disruptions and their impact on transport networks, highlighting studies that assess the vulnerability of critical infrastructure components such as roads, bridges, and tunnels in seismically active regions.

Further sections explore the integration of graph theory and GIS methodologies for resilience assessments, examining how these tools have been applied to urban and interregional transport networks. Studies that combine network topology with spatial data have proven effective in predicting disruption impacts and guiding infrastructure resilience strategies. The final section synthesizes findings from previous research and identifies gaps in the existing literature, particularly concerning the assessment of interregional transport networks, which often face unique challenges due to their scale, jurisdictional boundaries, and varied infrastructure conditions.

By reviewing these methodologies and case studies, this chapter establishes a foundation for analyzing transport network resilience using graph-based indicators in an interregional context. The findings from the literature will inform the methodological approach presented in the next chapter, ensuring a comprehensive framework for evaluating and improving transport network resilience across regions.

\section{Graph Theory Applications in Transport Network Resilience}
Graph theory has emerged as a powerful tool in the study of transport network resilience, providing a mathematical framework to analyze network connectivity, structure, and redundancy. By representing transport networks as graphs composed of nodes (intersections, stations) and edges (roads, railways), graph theory enables the evaluation of how networks function under normal conditions and how they respond to disruptions.

Historically, transport network vulnerability assessments have largely relied on accessibility-based approaches, where the impact of disruptions is measured in terms of accessibility loss and travel cost increases. Many studies have evaluated network vulnerability by assessing connectivity reliability, generalized travel cost, and accessibility indicators \cite{taylor2006, jenelius2006}. These methods focus on how disruptions affect the ease of reaching destinations and the efficiency of transport services. However, graph theory provides a complementary and more structured approach, allowing for the quantification of redundancy, network robustness, and interconnectivity, which are essential for understanding resilience.

In this context, \cite{thottolil2021} investigate the topology of Bangalore’s road networks, focusing on connectivity and coverage through graph theory. Their research highlights the relationship between road connectivity and the development of the road transport system, utilizing OpenStreetMap data to compute various topological measures. Through correlation analysis, they demonstrate the strength of the relationship among these indicators and propose a model linking connectivity indices with network density metrics. This model offers valuable insights into how urban transport networks can be optimized for efficiency and resilience.

Graph-based approaches in transport resilience studies have demonstrated their effectiveness in assessing the structural properties of networks, identifying critical links and nodes, and predicting the impact of disruptions on overall connectivity. The ability of graph theory to quantify network robustness allows for a systematic evaluation of different resilience indicators, including connectivity indices, redundancy measures, and centrality metrics. These indicators provide transport planners with the necessary tools to design networks that can withstand disruptions and maintain functional accessibility.

By applying graph theory to transport resilience analysis, researchers have been able to develop models that can be adapted to various urban and interregional contexts. These models assist in the optimization of infrastructure investments, ensuring that road networks remain efficient and reliable under different operational conditions. The ability to integrate graph-theoretic analysis with real-world transport data allows for more informed decision-making, contributing to the development of sustainable and resilient transportation systems.
\section{Seismic Disruptions and Their Impact on Transport Networks}
Seismic disruptions significantly affect the connectivity, redundancy, and resilience of transportation networks in seismically active regions. Earthquakes can cause severe structural damage to transport infrastructure, including roadways, bridges, tunnels, and railway systems, leading to prolonged accessibility disruptions and economic losses. The ability of a transport network to withstand and recover from such disruptions is critical for ensuring mobility, emergency response effectiveness, and overall regional stability.

\cite{dou2023} highlight the critical need for resilience assessments of road transportation infrastructure systems under earthquake impacts, emphasizing the importance of connectivity reliability and travel time evaluations for emergency planning. Their study on the Aba Autonomous Prefecture demonstrates how Monte Carlo simulations can effectively model uncertainties in seismic hazard assessments, revealing vulnerabilities in road assets such as bridges and tunnels. Their findings indicate that the failure of critical transport links during seismic events can lead to network fragmentation, severely affecting evacuation routes and access to essential services.

Seismic risk modeling has increasingly focused on integrating graph-based resilience metrics with geospatial data to assess transport network vulnerabilities. By combining redundancy indicators and structural failure probability models, researchers can identify high-risk zones and propose mitigation strategies to minimize seismic-induced disruptions. Understanding how network connectivity changes due to seismic damage is essential for designing adaptive transport systems that can reroute traffic efficiently when primary corridors become impassable.

The necessity of incorporating connectivity and redundancy metrics into resilience frameworks is underscored by studies such as Dou et al. (2023), which advocate for a data-driven approach to strengthening transport networks in earthquake-prone areas. Their research suggests that proactive infrastructure reinforcement, emergency detour planning, and strategic investment in redundant links can significantly enhance the post-disruption performance of road transportation systems. These findings align with broader efforts to develop seismic resilience strategies that ensure transport networks remain functional even in the aftermath of catastrophic events.

\section{Graph Theory and GIS-Based Resilience Assessment}
Assessing the resilience of transport networks under earthquake-induced disruptions is critical for enhancing community safety and economic stability. Recent studies have highlighted effective methodologies that leverage graph theory and Geographic Information Systems (GIS) to evaluate such resilience. By integrating network topology with spatial data, researchers can assess the vulnerability of transport infrastructure, identify critical network components, and develop targeted mitigation strategies to enhance overall resilience.

\cite{rouhana2020} introduced a novel approach based on graph theory concepts, focusing on the topological properties of the transportation network in Beirut. Their research demonstrated how the road network's configuration influences its resilience against random and targeted failures, revealing that critical network components could be identified to optimize resource allocation and minimize downtime. Their findings emphasize the importance of network redundancy and structural robustness in maintaining connectivity under seismic stress.

In a complementary study \cite{ dapuzzo2022} proposed an innovative methodological framework that integrates geo-statistical approaches to analyze soil properties and predict liquefaction risks. This framework, applied to a district in northern Italy, assessed the economic losses associated with seismic-induced road disruptions, providing a comprehensive evaluation of accessibility impacts during seismic events. The study highlights how integrating geological vulnerability assessments with network topology analysis allows for a more accurate prediction of post-disruption accessibility challenges.

Both studies underscore the importance of utilizing graph-based and GIS-based methodologies to quantify resilience effectively. Graph theory enables researchers to measure network connectivity, redundancy, and vulnerability, while GIS applications provide spatial insights into potential risk factors such as seismic zones, population density, and emergency service accessibility. By combining topological analysis with spatial data, researchers can develop robust planning tools that enhance the resilience of transport networks against earthquake-induced disruptions. These tools contribute to sustainable urban development and disaster preparedness by enabling decision-makers to identify weak points in transport infrastructure and prioritize strategic interventions to strengthen resilience.

\section{Graph-Based Indicators for Transport Network Resilience}
Graph-based indicators play a fundamental role in assessing the resilience of transport networks, particularly in evaluating their connectivity, redundancy, and ability to sustain functionality after disruptions. Various studies have explored how key graph theory metrics, including Beta Index, Gamma Index, Clustering Coefficient, and Betweenness Centrality, contribute to understanding transport network robustness. These indicators provide insights into network structure, redundancy, and critical vulnerability points, helping planners and researchers design more resilient transport systems.

 \cite{wu2018} investigated the resilience of transportation networks using graph theory measures. Their study demonstrated that the Beta Index and Gamma Index are effective in quantifying network connectivity and redundancy. Networks with higher Beta and Gamma values exhibited greater resilience to disruptions due to the availability of multiple alternative routes. The authors also emphasized the importance of Betweenness Centrality in identifying critical nodes whose failure could severely impact network performance.

\cite{leu2010}conducted a case study on Melbourne's ground transportation networks, where they utilized the Clustering Coefficient to assess local connectivity and redundancy. Their findings revealed that networks with higher clustering coefficients, particularly the tram network, demonstrated superior resilience to both random and targeted disruptions. They concluded that the Clustering Coefficient serves as a reliable indicator of a network's ability to maintain connectivity despite node or link failures.

\cite{ducruet2024} provided a comprehensive overview of graph theory measures for transportation networks, emphasizing that the Beta Index measures the level of connectivity, with values greater than 1 indicating a more complex network with higher redundancy. They described the Gamma Index as a metric for evaluating connectivity relative to the maximum possible connections, stating that networks with Gamma values closer to 1 exhibit greater redundancy and resilience to disruptions.

\cite{xu2024} examined the relationship between network structure and resilience, finding that networks with higher clustering coefficients performed better in resilience dimension reduction, indicating a stronger ability to withstand disruptions. They also noted that Betweenness Centrality is particularly useful for identifying critical nodes and potential vulnerabilities within the network structure.

\cite{hui2020} explored methods for improving the resilience of multilayer supply chain networks, which, despite their primary focus on supply chains, have implications for transport networks. The study employed Betweenness Centrality to identify critical nodes and proposed a layer-based rewiring algorithm to enhance network resilience against both random and targeted disruptions.

\cite{ngogo2017} discussed measures of line distributions in transport networks, highlighting that higher Beta Index values indicate better connectivity and integration within a network. The study suggested that transport networks with higher Beta Index values exhibit greater redundancy and are therefore more resilient to disruptions.

These studies collectively demonstrate that higher values of Beta Index, Gamma Index, and Clustering Coefficient correlate with greater network redundancy and resilience to disruptions. Betweenness Centrality emerges as a crucial metric for identifying critical nodes whose failure could significantly impact network performance. Together, these indicators provide valuable insights into a transport network's ability to maintain functionality in the face of disruptions, offering planners and policymakers essential tools for designing robust and resilient transport systems.
\section{Integration of Graph Theory and GIS for Seismic Resilience}

The integration of graph theory-based resilience metrics with GIS-driven spatial modeling presents a robust framework for assessing transport network disruptions, especially in regions prone to seismic activity. By combining topological analysis with spatial data, researchers can evaluate network vulnerabilities, predict disruption impacts, and develop effective mitigation strategies. This approach enhances the ability of transportation systems to withstand seismic shocks while ensuring connectivity and accessibility remain intact.

\cite{ganin2017} emphasize the resilience and efficiency of urban transportation networks, highlighting how disruptions can expose inefficiencies while also revealing the inherent resilience capabilities of these systems under stress conditions. Their research suggests that evaluating both efficiency and resilience metrics simultaneously provides a more comprehensive understanding of network behavior during seismic events. This dual consideration can be instrumental in predictive modeling frameworks, helping decision-makers design more adaptive transport networks.

\cite{xu2023} examine intermodal connectivity within public transportation systems in Hong Kong, illustrating how enhanced network integration can mitigate vulnerabilities associated with unforeseen disruptions such as earthquakes. Their findings indicate that improving intermodal links not only strengthens operational robustness but also facilitates faster recovery after disruptions. The use of GIS applications in their study provides valuable spatial data for optimizing transport system design and improving post-disaster response strategies.

\cite{zuo2022} present a case study focusing on New York City's transportation networks following severe weather events, utilizing graph theory to quantify network performance under adverse conditions. Their study identifies critical nodes whose failure could disproportionately affect overall system functionality, emphasizing the importance of integrating spatial analysis techniques to enhance predictive accuracy regarding potential disruption impacts from seismic activities. Their findings underscore the need for transport networks to incorporate redundancy and adaptive capacity to maintain service continuity during crises.

\cite{alipour2022} propose a comprehensive framework for assessing land transportation networks' resilience through graph theoretical approaches combined with optimization models. Their study accounts for node vulnerabilities and redundancy within the network structure, offering methodologies crucial for evaluating potential impacts from seismic disturbances. By integrating GIS with graph-based resilience assessments, their framework enables the identification of high-risk transport corridors and the development of strategic reinforcements.

\cite{cox2019} explore transit network resilience through empirical analysis using graph theory metrics specifically designed to evaluate performance during operational disruptions such as those caused by earthquakes. Their empirical approach provides actionable insights for improving infrastructure planning and risk mitigation strategies based on real-world data. Their study suggests that incorporating empirical resilience metrics into GIS-based spatial planning enhances infrastructure preparedness, reducing long-term economic and social impacts from transport network failures.

\cite{chen2016} contribute an innovative approach by employing GIS-based tools alongside big data analytics to model roadway and transit system resiliency concerning neighborhood-level characteristics in New York City post-hurricane scenarios. Their findings emphasize the significance of localized spatial patterns in understanding broader systemic vulnerabilities within transport networks exposed to seismic threats. By incorporating high-resolution spatial data into resilience modeling, their study offers valuable insights into adaptive transport planning at multiple scales.

These studies collectively affirm that combining graph theory-based metrics with GIS-driven spatial modeling creates a powerful predictive framework for assessing transport network disruptions in seismically active regions. The integration of topological network analysis with spatial risk modeling enhances decision-making capabilities by identifying weak points in the transport infrastructure, predicting potential failures, and developing targeted intervention strategies. This approach contributes to the design of more resilient transport networks that can better withstand seismic disturbances while maintaining operational efficiency and accessibility.

\section{Summary and Research Gaps}

This literature review has explored key methodologies and indices used to assess the resilience of transportation networks, with a particular emphasis on graph theory indicators and GIS-based spatial modeling. Studies have demonstrated that graph-based indices such as beta index, gamma index, clustering coefficient, and betweenness centrality play a fundamental role in quantifying connectivity, redundancy, and network robustness. These indicators allow for an evaluation of network performance under normal and disrupted conditions, helping researchers identify critical nodes, vulnerable links, and alternative routes. Furthermore, studies integrating GIS-based spatial analysis have highlighted the importance of incorporating real-world geographical constraints into network resilience modeling, particularly in regions exposed to seismic hazards.

Despite the advancements in transport resilience research, there is a lack of synthesis and standardization in methodological approaches. Different studies have applied varied graph-theoretical models, resilience metrics, and simulation techniques, often focusing on urban transit networks rather than regional and interregional transport systems. While some research has analyzed seismic disruptions, much of the existing work has been conducted on city-scale public transport systems, major road networks in metropolitan areas, or international multimodal corridors. These approaches may not be directly applicable to extra-urban rural roads, which present unique structural and functional characteristics.

A key limitation in previous studies is that they have primarily focused on general network connectivity and vulnerability assessments without a targeted analysis of interregional transport networks in rural and semi-rural settings. Many studies have concentrated on high-density urban areas, emphasizing the resilience of rail networks, metro systems, or urban road grids, while less attention has been given to extra-urban and rural road networks that are vital for interregional mobility. Additionally, while some research has integrated seismic risk analysis with network resilience assessments, there is still a lack of a unified framework that systematically evaluates redundancy and network robustness across different road classifications such as motorways, trunk roads, primary roads, and secondary roads in seismically active regions.

This study aims to address these gaps by providing a comprehensive resilience assessment of the extra-urban rural road network in Southern Italy, covering seven regions that are crucial for interregional mobility. Unlike previous studies that primarily focused on urban transport systems, this research evaluates road classes including motorways, trunk roads, primary roads, and secondary roads, which serve as key interregional connectors. Furthermore, this study integrates graph-based resilience metrics with GIS spatial modeling, applying them specifically to a seismically active region where network disruptions could have significant socio-economic consequences.

By incorporating graph theory indicators, connectivity metrics, and redundancy assessments, the present study seeks to bridge the gap between urban-focused resilience studies and interregional transport network analysis. It aims to develop a standardized methodology for evaluating transport network resilience across different road classifications, emphasizing seismic hazard impacts on extra-urban road infrastructure. The findings of this research will contribute to more effective transport planning and disaster preparedness strategies, ensuring that interregional road networks in Southern Italy are more resilient, adaptive, and capable of sustaining mobility during disruptive events.
